\begin{ex}
Số gồm 6 chục nghìn, 4 nghìn, 5 trăm và 2 đơn vị viết là gì?
\shortans{64502}
\loigiai{
Số đó là 64502.
}
\end{ex}

\begin{ex}
Số liền sau của số 89 999 là số nào?
\shortans{90000}
\loigiai{
$89999 + 1 = 90000$.
}
\end{ex}

\begin{ex}
Số liền trước của số 100 000 là số nào?
\shortans{99999}
\loigiai{
$100000 - 1 = 99999$.
}
\end{ex}

\begin{ex}
Tính: $35000 + 45000 = \dots$
\shortans{80000}
\loigiai{
$35000 + 45000 = 80000$.
}
\end{ex}

\begin{ex}
Tính: $90000 - 35000 = \dots$
\shortans{55000}
\loigiai{
$90000 - 35000 = 55000$.
}
\end{ex}

\begin{ex}
Tính: $16000 \times 4 = \dots$
\shortans{64000}
\loigiai{
$16000 \times 4 = 64000$.
}
\end{ex}

\begin{ex}
Tính: $84000 : 4 = \dots$
\shortans{21000}
\loigiai{
$84000 : 4 = 21000$.
}
\end{ex}

\begin{ex}
Tính giá trị biểu thức: $15000 \times 3 + 25000$.
\shortans{70000}
\loigiai{
$45000 + 25000 = 70000$.
}
\end{ex}

\begin{ex}
Tính giá trị biểu thức: $100000 - 60000 : 2$.
\shortans{70000}
\loigiai{
$100000 - 30000 = 70000$.
}
\end{ex}

\begin{ex}
Tìm $x$: $x \times 5 = 60000$.
\shortans{12000}
\loigiai{
$x = 60000 : 5 = 12000$.
}
\end{ex}

\begin{ex}
Tìm $x$: $x : 3 = 24000$.
\shortans{72000}
\loigiai{
$x = 24000 \times 3 = 72000$.
}
\end{ex}

\begin{ex}
Một hình chữ nhật có chiều dài 15 cm, chiều rộng 8 cm. Chu vi của nó là bao nhiêu xăng-ti-mét?
\shortans{46}
\loigiai{
$(15 + 8) \times 2 = 46$ cm.
}
\end{ex}

\begin{ex}
Một hình chữ nhật có chiều dài 15 cm, chiều rộng 8 cm. Diện tích của nó là bao nhiêu xăng-ti-mét vuông?
\shortans{120}
\loigiai{
$15 \times 8 = 120\text{ cm}^2$.
}
\end{ex}

\begin{ex}
Một hình vuông có chu vi 36 cm. Cạnh của hình vuông đó dài bao nhiêu xăng-ti-mét?
\shortans{9}
\loigiai{
$36 : 4 = 9$ cm.
}
\end{ex}

\begin{ex}
Một hình vuông có cạnh dài 9 cm. Diện tích của nó là bao nhiêu xăng-ti-mét vuông?
\shortans{81}
\loigiai{
$9 \times 9 = 81\text{ cm}^2$.
}
\end{ex}

\begin{ex}
Hình tròn có bán kính 6 cm thì đường kính là bao nhiêu xăng-ti-mét?
\shortans{12}
\loigiai{
$6 \times 2 = 12$ cm.
}
\end{ex}

\begin{ex}
Khối lập phương có bao nhiêu đỉnh?
\shortans{8}
\loigiai{
Có 8 đỉnh.
}
\end{ex}

\begin{ex}
Khối hộp chữ nhật có bao nhiêu mặt?
\shortans{6}
\loigiai{
Có 6 mặt.
}
\end{ex}

\begin{ex}
Đổi 1 tờ 100 000 đồng lấy các tờ 20 000 đồng được bao nhiêu tờ?
\shortans{5}
\loigiai{
$100000 : 20000 = 5$ tờ.
}
\end{ex}

\begin{ex}
Mua 4 chiếc bút bi giá 5 000 đồng mỗi chiếc. Đưa tờ 50 000 đồng thì được trả lại bao nhiêu đồng?
\shortans{30000}
\loigiai{
$50000 - 4 \times 5000 = 30000$ đồng.
}
\end{ex}

\begin{ex}
Tháng 8 có bao nhiêu ngày?
\shortans{31}
\loigiai{
Tháng 8 có 31 ngày.
}
\end{ex}

\begin{ex}
Tháng 9 có bao nhiêu ngày?
\shortans{30}
\loigiai{
Tháng 9 có 30 ngày.
}
\end{ex}

\begin{ex}
Một ngày có bao nhiêu giờ?
\shortans{24}
\loigiai{
Có 24 giờ.
}
\end{ex}

\begin{ex}
Làm tròn số 64 800 đến hàng nghìn ta được số nào?
\shortans{65000}
\loigiai{
Làm tròn lên 65 000.
}
\end{ex}

\begin{ex}
Số La Mã XV tương ứng với số tự nhiên nào?
\shortans{15}
\loigiai{
XV là 15.
}
\end{ex}
